\documentclass[10pt,a4paper]{amsart}
\usepackage[margin=19mm]{geometry}
\usepackage{amsmath,amssymb,mathtools}
\usepackage[scr=rsfs,cal=euler]{mathalfa}
\usepackage{xfrac}
\usepackage[unicode,hidelinks,bookmarks=false]{hyperref}
\newcommand{\ie}{i.e.\ }
\newcommand{\cf}{cf.\ }
\newcommand{\resp}{respectively\ }
\newcommand{\Subsection}{\subsection}
\providecommand{\nobreakdash}{\nobreak\hbox{-}}
\setlength{\emergencystretch}{2em}
\multlinegap0pt
\usepackage{amssymb}
\usepackage{algorithm}\usepackage{algpseudocode}
\usepackage{tikz}
\newtheorem{lemma}{Lemma}
\newcommand{\cut}{\mathop{\mathrm{cut}}}
\title{The antiferromagnetic Ising model\\ beyond line graphs}
\author{Mark Jerrum}
\date{Source: arXiv:2602.14915v1 (CC BY 4.0)}
\begin{document}
\maketitle

\noindent\emph{Source and licence.} The antiferromagnetic Ising model\\ beyond line graphs. Version arXiv:2602.14915v1 (CC BY 4.0), distributed by arXiv under the Creative Commons Attribution 4.0 International licence, http://creativecommons.org/licenses/by/4.0/, as recorded in the arXiv licence metadata for this exact version and shown on the abstract page https://arxiv.org/abs/2602.14915v1. Full licence notice: \texttt{licenses/arxiv-2602.14915-CC-BY-4.0.txt}.\par\smallskip

\noindent\textit{Editorial scope.} Counted unit: the article's lemma lem:Cgen (source0.tex lines 127--129). The article's complete original proof of it is delivered with it (source0.tex lines 131--133, one \texttt{proof} environment of 3 lines). No mathematics is written, repaired or re-derived: the statement and the proof are the article's own lines, in the article's own notation, and no retained line is substituted or re-flowed.  No further source-local context is needed: every cross-reference of every retained line resolves inside the two delivered spans.  Nothing was omitted from any retained span, so the package carries no omission record.  No retained line contains a citation, so the wrapper carries no bibliography and no bibliography entry.  No retained line contains a figure environment, an image-inclusion command or drawing code, so no image is delivered and the package declares no asset dependency.  Editorial formatting only, in addition: an AMS \texttt{amsart} preamble that restates the article's own declarations and macros used by the retained lines, namely the environments \texttt{lemma} on separate counters, together with the standard packages those lines need.  The retained environments print in this document as Lemma 1 in order of appearance, whereas the article prints the same items with the numbering of its own full document.  The complete native source of the article as distributed in the arXiv e-print for this exact version is bundled unchanged as \texttt{sources/194758/source0.tex}.
\begin{lemma}\label{lem:Cgen}
Every cubic graph $G$ has a cut of size at least $n=|V(G)|$.
\end{lemma}

\begin{proof}
Let $\sigma:V(G)\to\{+,-\}$ be a maximum cardinality cut in $G$.  If some vertex $v$ with $\sigma(v)={+}$ is adjacent to at least two vertices $u_1$, $u_2$ with $\sigma(u_1)=\sigma(u_2)={+}$ then we could obtain a larger cut by flipping the sign of~$v$.  It is similar when $\sigma(v)={-}$.  So every vertex in $G$ is adjacent to at least two vertices on the opposite side of the cut, and hence $|\cut\sigma|\geq n$.
\end{proof}

\end{document}
